\documentclass[../../main.tex]{subfiles}
\begin{document}

\begin{center}
  {\fontsize{48pt}{58pt}\selectfont \textbf{1976}\normalsize\textbf{年度}}
\end{center}

\begin{table}[htb]
  \centering
  \caption{出題テーマと難易度}
  \begin{tabular}{|c|p{9cm}|c|}
    \hline
    \textbf{問題番号} & \textbf{テーマ} & \textbf{難易度} \\
    \hline
    第1問 & 回転体の体積の最大化 & \\
    \hline
    第2問 & 直角に左折する動点の到達範囲 & \\
    \hline
    第3問 & 新課程：対称移動と回転／旧課程：円への接線と軌跡 & \\
    \hline
    第4問 & 有理関数の最小値を与える点の軌跡 & \\
    \hline
    第5問 & 新課程：円柱面の面積／旧課程：面積の変化率 & \\
    \hline
    第6問 & 新課程：行列の積と逆行列／旧課程：四次方程式の実根 & \\
    \hline
  \end{tabular}
\end{table}

\begin{table}[htb]
  \centering
  \caption{解答}
  \renewcommand{\arraystretch}{1.4}
  \begin{tabular}{|c|c|p{9cm}|}
    \hline
    \textbf{問題番号} & \textbf{小問} & \textbf{答え} \\
    \hline
    第1問 & -- & $0<c\le1$なら$(r,l)=(\sqrt{c},0)$．$c\ge1$なら$r=\dfrac{1+\sqrt{1+24c}}{6}$，$l=\dfrac{6c-1-\sqrt{1+24c}}{18}$ \\
    \hline
    第2問 & -- & 到達範囲は解答の図を参照 \\
    \hline
    第3問 & 新課程 & $f(r)=\sqrt{2}|r-2\sqrt{53}|$，$g(r)=\sqrt{2}(r+2\sqrt{53})$．$f(r)=0$となるのは$r=2\sqrt{53}$ \\
    \cline{2-3}
    & 旧課程 & 軌跡は解答の図を参照．長さは$5\sqrt{3}\pi$ \\
    \hline
    第4問 & (1) & $\alpha=\sqrt{1+t^2}-t$，$f(\alpha)=\dfrac{1}{\alpha^2}$（軌跡は解答参照） \\
    \cline{2-3}
    & (2) & $\beta=\min\{t,\alpha\}$．軌跡は解答の図を参照 \\
    \hline
    第5問 & 新課程 & 面積は$2\sqrt{3}+\dfrac{4\pi}{3}$ \\
    \cline{2-3}
    & 旧課程 & $S=\dfrac{1}{2}e^{2t}-2e^t+2t+\dfrac{3}{2}$ \\
    \hline
    第6問 & 新課程(i) & $R$は積について閉じている（証明問題） \\
    \cline{2-3}
    & 新課程(ii) & $a<-\dfrac{b^2}{4}$（図示は解答参照） \\
    \cline{2-3}
    & 旧課程(i) & 実根はすべての実数$k$で$4$個 \\
    \cline{2-3}
    & 旧課程(ii) & 曲線はある鉛直な直線に関して対称（証明問題） \\
    \hline
  \end{tabular}
\end{table}

\paragraph{試験形態}
6問．第3・5・6問は新旧課程別．試験時間は要確認．

\end{document}
